\chapter{Lógica Proposicional}

\section{Introdução}

A Lógica Proposicional constitui um dos fundamentos da Matemática e da Ciência da Computação. Seu objetivo é estudar proposições, seus valores lógicos e as formas pelas quais podem ser combinadas para construir argumentos válidos.

\section{Proposições}

\subsection{Definição}
 
Uma \textbf{proposição} é todo conjunto de palavras ou símbolos que exprime um pensamento de sentido completo, ao qual se pode atribuir um valor lógico: verdadeiro ou falso.

\subsection{Características}

As proposições possuem as seguintes características:

\begin{itemize}
    \item Declaram fatos;
    \item Exprimem juízos.
\end{itemize}

\subsection{Exemplos}

\begin{itemize}
    \item A Lua é um satélite da Terra.
    \item Recife é a capital de Pernambuco.
    \item $\pi > 5$
    \item $\sin\left(\frac{\pi}{2}\right) = 1$
\end{itemize}

\subsection{Não-Proposições}

Nem toda sentença é uma proposição. Perguntas, ordens ou exclamações não possuem valor lógico.

\begin{itemize}
    \item Qual horário o ônibus passará hoje?
\end{itemize}

\section{Princípios Fundamentais da Lógica}

\subsection{Princípio da Não-Contradição}

Uma proposição não pode ser verdadeira e falsa ao mesmo tempo.

\subsection{Princípio do Terceiro Excluído}

Toda proposição é ou verdadeira ou falsa, não existindo uma terceira possibilidade.

\subsection{Consequência}

Esses princípios implicam que a lógica clássica é \textbf{bivalente}, ou seja, admite apenas dois valores lógicos.

\subsection{Exemplos}

\begin{itemize}
    \item A Lua é um satélite da Terra. \hfill (Verdadeira)
    \item Vasco da Gama descobriu o Brasil. \hfill (Falsa)
    \item Recife é a capital de Pernambuco. \hfill (Verdadeira)
    \item Dante escreveu ``Os Lusíadas''. \hfill (Falsa)
\end{itemize}

\section{Tipos de Proposições}

\subsection{Proposições Simples}

Uma \textbf{proposição simples} (ou atômica) é aquela que não contém nenhuma outra proposição como parte.

\paragraph{Representação}

São geralmente representadas por letras minúsculas:
\[
p, q, r, \ldots
\]

\paragraph{Exemplos}

\begin{itemize}
    \item $p$: Carlos é careca.
    \item $q$: Pedro é estudante.
    \item $r$: O número 25 é um quadrado perfeito.
\end{itemize}

\subsection{Proposições Compostas}

Uma \textbf{proposição composta} (ou molecular) é formada pela combinação de duas ou mais proposições simples.

\paragraph{Representação}

Podem ser representadas por letras maiúsculas:
\[
P, Q, R, \ldots
\]

\paragraph{Exemplos}

\begin{itemize}
    \item $P$: Carlos é careca \textbf{e} Pedro é estudante.
    \item $Q$: Carlos é careca \textbf{ou} Pedro é estudante.
    \item $R$: \textbf{Se} Carlos é careca, \textbf{então} é infeliz.
\end{itemize}

Essas proposições também são chamadas de \textbf{fórmulas proposicionais}.

\section{Conectivos Lógicos}

\subsection{Definição}

Os conectivos são palavras ou símbolos utilizados para formar novas proposições a partir de outras.

\subsection{Principais Conectivos}

\begin{itemize}
    \item Conjunção: $p \land q$ (``e'')
    \item Disjunção: $p \lor q$ (``ou'')
    \item Negação: $\neg p$ (``não'')
    \item Condicional: $p \rightarrow q$ (``se... então'')
    \item Bicondicional: $p \leftrightarrow q$ (``se, e somente se'')
\end{itemize}

\subsection{Exemplos}

\begin{itemize}
    \item O número 6 é par \textbf{e} o número 8 é cubo perfeito.
    \item \textbf{Não} está chovendo.
    \item \textbf{Se} Jorge é engenheiro, \textbf{então} sabe Matemática.
    \item Um triângulo é equilátero \textbf{se, e somente se,} é equiângulo.
\end{itemize}

\section{Valor Lógico}

\subsection{Definição}

O valor lógico de uma proposição indica se ela é verdadeira ou falsa.

\subsection{Notação}

\[
V(p)
\]

\begin{itemize}
    \item Se $p$ é verdadeira, então $V(p) = V$;
    \item Se $p$ é falsa, então $V(p) = F$.
\end{itemize}

\subsection{Exemplos}

\begin{itemize}
    \item $p$: O Sol é verde. \hfill $V(p) = F$
    \item $q$: Um hexágono tem 9 diagonais. \hfill $V(q) = V$
    \item $r$: $2$ é raiz de $x^2 + 3x - 4 = 0$. \hfill $V(r) = F$
\end{itemize}

\section{Tabelas-Verdade}

\subsection{Princípio Fundamental}

O valor lógico de uma proposição composta depende exclusivamente dos valores lógicos das proposições simples que a compõem.

\subsection{Definição}

Uma \textbf{tabela-verdade} é um dispositivo que apresenta todas as possíveis combinações de valores lógicos de uma proposição.

\subsection{Número de Linhas}

Se uma proposição possui $n$ variáveis, então sua tabela-verdade terá:

\[
2^n \text{ linhas}
\]

\subsection{Exemplos}

\subsubsection{Uma Proposição}

\begin{center}
\begin{tabular}{|c|}
\hline
$p$ \\
\hline
V \\
F \\
\hline
\end{tabular}
\end{center}

\subsubsection{Duas Proposições}

\begin{center}
\begin{tabular}{|c|c|}
\hline
$p$ & $q$ \\
\hline
V & V \\
V & F \\
F & V \\
F & F \\
\hline
\end{tabular}
\end{center}

\subsubsection{Três Proposições}

\begin{center}
\begin{tabular}{|c|c|c|}
\hline
$p$ & $q$ & $r$ \\
\hline
V & V & V \\
V & V & F \\
V & F & V \\
V & F & F \\
F & V & V \\
F & V & F \\
F & F & V \\
F & F & F \\
\hline
\end{tabular}
\end{center}

\section{Considerações Finais}

Neste capítulo foram apresentados os conceitos fundamentais da lógica proposicional, incluindo a definição de proposições, seus tipos, conectivos lógicos e o uso de tabelas-verdade. Esses conceitos são essenciais para o desenvolvimento do raciocínio lógico formal e servirão de base para estudos mais avançados, como equivalências lógicas, inferência e lógica de predicados.